\section{Conclusion and Future Work}\label{sec:conclusion}

Summarizing our evaluation, we highlight the following directions for further investigation.

\begin{enumerate}
    \item \textbf{Memory consumption.} Interaction nets programming requires duplicator agents for data reuse. Inpla eagerly copies entire subnets during duplication, incurring substantial memory overhead relative to the F\# prototype. Vine adopts a different strategy but still demands significant memory. Can coding patterns, evaluation strategies, and memory management mechanisms be co-designed to substantially reduce memory consumption?
    
    \item \textbf{Alternative matrix representations.} Tree-based representations looks suboptimal: the recursive divide-and-conquer pattern induces aggressive duplication, traversal dominates cell handling, and performance is highly sensitive to reordering. Could coordinate list (COO) formats~--- operations on which rely on sorting, a task that parallelizes well in interaction nets~--- offer better memory efficiency and performance?
    
   \item \textbf{Load balancing and scaling.} The current configuration of Inpla scales poorly on small graphs but well on large ones. Vine generally scales better than Inpla. Scaling and performance of both tools remain sensitive to matrix reordering. Can the load balancing strategy be improved to address these issues?
    
    \item \textbf{Advanced fusion techniques.} Optimization techniques from functional programming, such as supercompilation~\cite{10.1007/3-540-62064-8_20} or distillation~\cite{10.1007/978-3-030-98869-2_7}, could automate the fusion of operation sequences, analogous to manual kernel fusion for masks and other constructs in SuiteSparse:GraphBLAS. While partial evaluation has been shown possible for interaction nets~\cite{DBLP:conf/sas/Bechet92}, whether these more advanced techniques are adaptable remains an open question.
    
    \item \textbf{FPGA acceleration.} Could an FPGA designed specifically for interaction net reduction exploit the model's fine-grained parallelism, overcome scaling limitations of general-purpose hardware, and achieve substantial performance gains?
\end{enumerate}

Finally, we believe our work motivates comparisons among different interaction net implementations~--- in particular, by extending our evaluation to other systems such as HVM~--- and helps make interaction nets applicable to real-world problems.